% !TEX program = pdflatex
% ============================================================================
%  Unidade 5 - Gerenciando dados textuais com pandas
%  Trabalhando com Texto em Python I  -  CiberExt 26-29 / FEELT38103 / UFU
%
%  COMPILAR COM:  pdflatex unidade5_gerenciando_dados_pandas.tex   (rodar 2x)
%
%  Pacotes necessarios (todos presentes no MacTeX / TeX Live completo):
%    babel, inputenc, fontenc, lmodern, helvet, xcolor, geometry, listings,
%    tcolorbox, titlesec, enumitem, hyperref, microtype, booktabs, colortbl, array
% ============================================================================
\documentclass[11pt,a4paper]{article}

\usepackage[utf8]{inputenc}
\usepackage[T1]{fontenc}
\usepackage[brazilian]{babel}
\usepackage{lmodern}
\usepackage[scaled=0.92]{helvet}
\renewcommand{\familydefault}{\sfdefault}
\usepackage{microtype}
\usepackage[a4paper,margin=2.4cm]{geometry}
\usepackage{xcolor}
\usepackage{enumitem}
\usepackage{array}
\usepackage{listings}
\usepackage[most]{tcolorbox}
\tcbuselibrary{listings,skins,breakable}
\usepackage{titlesec}
\usepackage{booktabs}
\usepackage{colortbl}
\usepackage{hyperref}

% -------------------------------------------------------------- paleta ------
\definecolor{accent}{HTML}{6C4CF1}
\definecolor{cyan}{HTML}{00B4D8}
\definecolor{subink}{HTML}{2B3242}
\definecolor{codebg}{HTML}{1E2430}
\definecolor{codetext}{HTML}{E6E9EF}
\definecolor{ccomment}{HTML}{7FD18B}
\definecolor{cstring}{HTML}{F0B866}
\definecolor{ckw}{HTML}{C792EA}
\definecolor{cbuiltin}{HTML}{82AAFF}
\definecolor{outbg}{HTML}{F0F2F7}
\definecolor{outtext}{HTML}{2F3646}
\definecolor{outframe}{HTML}{9AA4B8}
\definecolor{notec}{HTML}{3B82F6}
\definecolor{tipc}{HTML}{10B981}
\definecolor{warnc}{HTML}{F59E0B}
\definecolor{inlink}{HTML}{5A34D6}
\definecolor{headrow}{HTML}{EFE9FF}

\hypersetup{colorlinks=true,linkcolor=accent,urlcolor=cyan,citecolor=accent,
  pdftitle={Gerenciando dados textuais com pandas},pdfauthor={CiberExt 26-29 / FEELT38103 / UFU}}

% --------------------------------------------------- estilo do codigo -------
\lstset{
  extendedchars=true,
  breaklines=true,
  breakatwhitespace=false,
  keepspaces=true,
  columns=fullflexible,
  tabsize=4,
  showstringspaces=false,
  literate=%
    {á}{{\'a}}1 {é}{{\'e}}1 {í}{{\'i}}1 {ó}{{\'o}}1 {ú}{{\'u}}1
    {â}{{\^a}}1 {ê}{{\^e}}1 {ô}{{\^o}}1
    {ã}{{\~a}}1 {õ}{{\~o}}1 {à}{{\`a}}1 {ü}{{\"u}}1 {ä}{{\"a}}1 {ö}{{\"o}}1
    {ç}{{\c c}}1 {Ç}{{\c C}}1
    {Á}{{\'A}}1 {É}{{\'E}}1 {Í}{{\'I}}1 {Ó}{{\'O}}1 {Ú}{{\'U}}1
    {Ã}{{\~A}}1 {Õ}{{\~O}}1 {Â}{{\^A}}1 {Ê}{{\^E}}1 {À}{{\`A}}1
    {—}{{---}}1 {–}{{--}}1
    {“}{{``}}1 {”}{{''}}1 {‘}{{`}}1 {’}{{'}}1 {´}{{'}}1
}

\lstdefinestyle{pystyle}{
  language=Python,
  basicstyle=\ttfamily\small\color{codetext},
  keywordstyle=\color{ckw}\bfseries,
  keywordstyle=[2]\color{cbuiltin},
  morekeywords=[2]{print,type,len,range,enumerate,list,append,extend,read_csv,DataFrame,apply,spacy,pd},
  commentstyle=\color{ccomment}\itshape,
  stringstyle=\color{cstring},
}
\lstdefinestyle{outstyle}{
  basicstyle=\ttfamily\small\color{outtext},
}

% caixa de codigo Python (fundo escuro, barra lateral colorida + titulo)
\newtcblisting{pybox}{
  breakable, enhanced,
  listing only, listing options={style=pystyle},
  colback=codebg, colframe=codebg, boxrule=0pt, arc=3pt,
  borderline west={3pt}{0pt}{accent},
  left=10pt,right=10pt,top=6pt,bottom=6pt,
  title=\textbf{Python}, coltitle=white, colbacktitle=accent,
  fonttitle=\sffamily\bfseries\scriptsize,
  before skip=10pt, after skip=12pt,
}
% caixa de saida (fundo claro)
\newtcblisting{outbox}{
  breakable, enhanced,
  listing only, listing options={style=outstyle},
  colback=outbg, colframe=outframe, boxrule=0.4pt, arc=3pt,
  left=10pt,right=10pt,top=6pt,bottom=6pt,
  title=\textbf{Saída}, coltitle=white, colbacktitle=outframe,
  fonttitle=\sffamily\bfseries\scriptsize,
  before skip=10pt, after skip=12pt,
}

% ------------------------------------------------------ caixas de aviso -----
\newtcolorbox{notebox}[1][Nota]{enhanced,breakable,
  colback=notec!6,colframe=notec,coltitle=white,title=#1,
  fonttitle=\sffamily\bfseries\small,boxrule=0pt,leftrule=4pt,arc=3pt,
  left=10pt,right=10pt,top=7pt,bottom=7pt,before skip=10pt,after skip=12pt}
\newtcolorbox{tipbox}[1][Dica]{enhanced,breakable,
  colback=tipc!7,colframe=tipc,coltitle=white,title=#1,
  fonttitle=\sffamily\bfseries\small,boxrule=0pt,leftrule=4pt,arc=3pt,
  left=10pt,right=10pt,top=7pt,bottom=7pt,before skip=10pt,after skip=12pt}
\newtcolorbox{warnbox}[1][Atenção]{enhanced,breakable,
  colback=warnc!8,colframe=warnc,coltitle=white,title=#1,
  fonttitle=\sffamily\bfseries\small,boxrule=0pt,leftrule=4pt,arc=3pt,
  left=10pt,right=10pt,top=7pt,bottom=7pt,before skip=10pt,after skip=12pt}
\newtcolorbox{objbox}[1][Objetivos de aprendizagem]{enhanced,breakable,
  colback=accent!5,colframe=accent,coltitle=white,title=#1,
  fonttitle=\sffamily\bfseries\small,boxrule=0pt,leftrule=4pt,arc=3pt,
  left=10pt,right=10pt,top=7pt,bottom=8pt,before skip=10pt,after skip=14pt}

% ------------------------------------------------------ titulos de secao ----
\titleformat{\section}{\Large\bfseries\sffamily\color{accent}}
  {\thesection.}{0.5em}{}
\titleformat{\subsection}{\large\bfseries\sffamily\color{subink}}
  {\thesubsection}{0.5em}{}
\titlespacing*{\section}{0pt}{18pt}{6pt}
\titlespacing*{\subsection}{0pt}{12pt}{4pt}

% codigo em linha (verbatim, colorido) via atalho @...@
\lstMakeShortInline[basicstyle=\ttfamily\color{inlink}]{@}

\setlength{\parindent}{0pt}
\setlength{\parskip}{0.55em}

\newcommand{\thd}[1]{\textbf{#1}}

\begin{document}

% --------------------------------------------------------------- capa -------
\begin{tcolorbox}[enhanced, left color=accent, right color=cyan, coltext=white,
  boxrule=0pt, arc=8pt, left=20pt,right=20pt,top=16pt,bottom=18pt,
  before skip=0pt, after skip=20pt]
  {\footnotesize\bfseries CIBEREXT 26-29 \ \textbullet\ \ FEELT38103 \ \textbullet\ \ UNIVERSIDADE FEDERAL DE UBERLÂNDIA}\\[10pt]
  {\Huge\bfseries Gerenciando dados textuais com pandas}\\[8pt]
  {\large Trabalhando com Texto em Python I \ --- \ Unidade 5}\\[6pt]
  {\normalsize Importar, examinar, estender e salvar DataFrames; aplicar PLN a colunas.}\\[10pt]
  {\footnotesize Adaptado de \textit{Applied Language Technology} (Universidade de Helsinque) --- applied-language-technology.mooc.fi}
\end{tcolorbox}

\begin{objbox}
Ao final desta unidade, você deverá saber:
\begin{itemize}[leftmargin=1.3em,itemsep=2pt,topsep=3pt]
  \item importar dados para um @DataFrame@ do pandas;
  \item explorar dados armazenados em um @DataFrame@;
  \item acrescentar (estender) dados a um @DataFrame@;
  \item salvar os dados de um @DataFrame@.
\end{itemize}
\end{objbox}

O \textbf{pandas} é uma biblioteca do Python para trabalhar com \textbf{dados tabulares}. Começamos importando-a --- usamos @as@ para controlar o nome do módulo; por convenção, o pandas é abreviado como @pd@:

\begin{pybox}
import pandas as pd
\end{pybox}

% =========================================================== 1 ==============
\section{Importando dados para o pandas}

\subsection{A partir de um único arquivo}
Formatos típicos para distribuir corpora incluem \textbf{CSV} (\textit{Comma-separated Values}), \textbf{JSON} e texto simples. O pandas oferece muitas funções de leitura --- dá até para importar planilhas do Excel! O exemplo carrega um recorte do \textit{SFU Opinion and Comments Corpus} (SOCC), com artigos de opinião do jornal canadense \textit{The Globe and Mail}. A função @read_csv()@ recebe o \textbf{caminho} do arquivo:

\begin{pybox}
# Le o arquivo CSV e atribui a saida a variavel 'socc'
socc = pd.read_csv('data/socc_gnm_articles.csv')

# Examina o tipo do objeto guardado em 'socc'
type(socc)
\end{pybox}
\begin{outbox}
pandas.core.frame.DataFrame
\end{outbox}

O pandas devolve o conteúdo do CSV em um \textbf{@DataFrame@}, sua estrutura de dados nativa. O método @head()@ mostra as primeiras linhas:

\begin{pybox}
# Imprime as primeiras cinco linhas do DataFrame
socc.head(5)
\end{pybox}

O @DataFrame@ tem forma \textbf{tabular}, com colunas (@article_id@, @title@, @author@, @article_text@\ldots) e um \textbf{índice} por linha (0, 1, 2, \ldots). Versão compacta (algumas colunas omitidas para caber):

\begin{center}\small
\begin{tabular}{c c p{5.2cm} p{3cm} c}
\toprule
\rowcolor{headrow}\thd{índice} & \thd{article\_id} & \thd{title} & \thd{author} & \thd{ntop\_level\_comments} \\
\midrule
0 & 26842506 & The Tories deserve another mandate\ldots & GLOBE EDITORIAL & 1378.0 \\
1 & 26055892 & Harper hysteria a sign of closed liberal minds & Konrad Yakabuski & 455.0 \\
2 & 6929035 & Too many first nations people live in a dream\ldots & Jeffrey Simpson & 433.0 \\
3 & 19047636 & The Globe's editorial board endorses Tim Hudak\ldots & GLOBE EDITORIAL & 432.0 \\
4 & 11672346 & Disgruntled Arab states look to strip Canada\ldots & Campbell Clark & 411.0 \\
\bottomrule
\end{tabular}
\end{center}

O acessador @.at[]@ inspeciona \textbf{um único item}. O valor da coluna @title@ no índice 123:

\begin{pybox}
socc.at[123, 'title']
\end{pybox}
\begin{outbox}
"How Toronto got a 'world-class,' gold-plated, half-billion-dollar empty train"
\end{outbox}

\subsection{A partir de múltiplos arquivos}
Outro cenário comum é ter \textbf{vários arquivos} de texto. Primeiro, coletamos os arquivos com a classe @Path@ (vista na Unidade 1):

\begin{pybox}
# Importa a classe Path
from pathlib import Path

# Cria um objeto Path que aponta para o diretorio com os dados
corpus_dir = Path('data')

# Coleta todos os arquivos .txt no diretorio do corpus
corpus_files = list(corpus_dir.glob('*.txt'))

# Verifica os arquivos do corpus
corpus_files
\end{pybox}
\begin{outbox}
[PosixPath('data/WP_1990-08-10-25A.txt'),
 PosixPath('data/NYT_1991-01-16-A15.txt'),
 PosixPath('data/WP_1991-01-17-A1B.txt')]
\end{outbox}

Criamos um @DataFrame@ \textbf{vazio} definindo a forma de antemão: o número de linhas (@index@, com @range()@ de 0 até @len(corpus_files)@) e os nomes das colunas (@columns@):

\begin{pybox}
# Cria um DataFrame e atribui o resultado a variavel 'df'
df = pd.DataFrame(index=range(0, len(corpus_files)), columns=['filename', 'text'])

# Chama a variavel para inspecionar a saida
df
\end{pybox}

\begin{center}\small
\begin{tabular}{c l l}
\toprule
\rowcolor{headrow}\thd{índice} & \thd{filename} & \thd{text} \\
\midrule
0 & NaN & NaN \\
1 & NaN & NaN \\
2 & NaN & NaN \\
\bottomrule
\end{tabular}
\end{center}

Percorremos os @Path@s, lemos o conteúdo e o adicionamos com o acessador @.at@:

\begin{pybox}
# Percorre os arquivos do corpus e conta cada volta com enumerate()
for i, f in enumerate(corpus_files):

    # Le o conteudo do arquivo
    text = f.read_text(encoding='utf-8')

    # Pega o nome do arquivo do objeto Path
    filename = f.name

    # Atribui o texto do arquivo ao indice 'i' na coluna 'text' usando
    # o acessador .at -- isso modifica o DataFrame "no lugar" (in place);
    # nao e preciso atribuir o resultado a uma variavel
    df.at[i, 'text'] = text

    # Faz o mesmo com o nome do arquivo
    df.at[i, 'filename'] = filename
\end{pybox}

O @DataFrame@ fica populado com os nomes dos arquivos e o texto:

\begin{center}\small
\begin{tabular}{c l p{7cm}}
\toprule
\rowcolor{headrow}\thd{índice} & \thd{filename} & \thd{text} \\
\midrule
0 & WP\_1990-08-10-25A.txt & *We Don't Stand for Bullies': Diverse Voices\ldots \\
1 & NYT\_1991-01-16-A15.txt & U.S. TAKING STEPS TO CURB TERRORISM: F.B.I. I\ldots \\
2 & WP\_1991-01-17-A1B.txt & U.S., Allies Launch Massive Air War Against T\ldots \\
\bottomrule
\end{tabular}
\end{center}

% =========================================================== 2 ==============
\section{Examinando DataFrames}

Os @DataFrame@s guardam muita informação, organizada em \textbf{colunas}, acessíveis pelo atributo @columns@:

\begin{pybox}
# Recupera as colunas e seus nomes
socc.columns
\end{pybox}
\begin{outbox}
Index(['article_id', 'title', 'article_url', 'author', 'published_date',
       'ncomments', 'ntop_level_comments', 'article_text'],
      dtype='object')
\end{outbox}

Uma coluna inteira é acessada com colchetes @[]@ e o nome como \textit{string} --- como as chaves de um dicionário:

\begin{pybox}
# Recupera o conteudo da coluna 'author' no DataFrame 'socc'
socc['author']
\end{pybox}
\begin{outbox}
0         GLOBE EDITORIAL
1        Konrad Yakabuski
2         Jeffrey Simpson
             ...
10337      Adam Radwanski
10338     GLOBE EDITORIAL
Name: author, Length: 10339, dtype: object
\end{outbox}

Cada coluna é, na verdade, um objeto \textbf{@Series@} --- pense no @DataFrame@ como a tabela inteira, cujas colunas são @Series@:

\begin{pybox}
# Verifica o tipo de 'socc' e de 'socc['author']'
type(socc), type(socc['author'])
\end{pybox}
\begin{outbox}
(pandas.core.frame.DataFrame, pandas.core.series.Series)
\end{outbox}

\begin{notebox}
Ao imprimir um @DataFrame@ ou @Series@, o pandas \textbf{omite} tudo entre as cinco primeiras e as cinco últimas linhas por padrão --- conveniente com milhares de linhas.
\end{notebox}

O método @value_counts()@ conta os valores únicos de uma @Series@:

\begin{pybox}
# Conta os valores unicos na coluna 'author'
socc['author'].value_counts()
\end{pybox}
\begin{outbox}
GLOBE EDITORIAL                   2712
Jeffrey Simpson                    649
Margaret Wente                     547
                                  ...
Jessica Scott-Reid                   1
Kenneth Oppel                        1
Name: author, Length: 1896, dtype: int64
\end{outbox}

Podemos \textbf{visualizar} o resultado com o método @.plot()@, que usa a biblioteca \textbf{matplotlib}. Passamos @bar@ ao argumento @kind@ e @[:10]@ para os dez autores mais prolíficos:

\begin{pybox}
# Magia do Jupyter que permite renderizar graficos do matplotlib no notebook!
# Voce so precisa executar este comando uma vez.
%matplotlib inline

# Conta os valores da coluna 'author' e limita ao top-10 antes de plotar.
socc['author'].value_counts()[:10].plot(kind='bar')
\end{pybox}

\begin{tipbox}
O @.plot()@ gera uma \textbf{imagem} (gráfico de barras) exibida no notebook; a saída de texto é apenas @<AxesSubplot: >@. Para outros tipos, mude @kind@ (ex.: @'line'@, @'hist'@, @'pie'@).
\end{tipbox}

Para colunas \textbf{numéricas}, o método @describe()@ dá estatísticas descritivas:

\begin{pybox}
# Estatisticas descritivas basicas da coluna 'ntop_level_comments'
socc['ntop_level_comments'].describe()
\end{pybox}
\begin{outbox}
count    10339.000000
mean        26.384273
std         39.786923
min          0.000000
25%          1.000000
50%         14.000000
75%         35.000000
max       1378.000000
Name: ntop_level_comments, dtype: float64
\end{outbox}

Lendo a saída: 10\,339 linhas; média (@mean@) $\approx 26$ comentários, com forte variação (desvio-padrão @std@ $\approx 40$); mínimo (@min@) 0; primeiro quartil (@25%@) $= 1$; mediana (@50%@) $= 14$; terceiro quartil (@75%@) $= 35$; máximo (@max@) $= 1\,378$.

\subsection{Filtrando linhas com \textnormal{@.loc@}}
Como achar os artigos com \textbf{zero} comentários? Usamos @.loc@ para selecionar linhas por valor. Como @=@ é reservado para atribuição, a comparação ``é igual a'' usa \textbf{dois} sinais (@==@):

\begin{pybox}
# Pega as linhas sem comentarios de nivel superior
socc.loc[socc['ntop_level_comments'] == 0]
\end{pybox}
\begin{outbox}
[... 2542 rows x 8 columns ...]
\end{outbox}

Retorna 2\,542 linhas. Para visões mais complexas, combinamos critérios com @&@ (``E'' lógico) --- cada critério entre \textbf{parênteses} @()@:

\begin{pybox}
# Numero de comentarios de nivel superior para o autor Hayden King
socc.loc[(socc['ntop_level_comments'] == 0) & (socc['author'] == 'Hayden King')]
\end{pybox}

Retorna apenas a linha de índice 7797 --- o único artigo dele sem comentários.

% =========================================================== 3 ==============
\section{Estendendo DataFrames}

É fácil \textbf{acrescentar} informação. Um cenário comum: carregar dados, analisar e guardar os resultados no mesmo @DataFrame@. Adicionamos uma coluna vazia com @[]@ e o tipo @None@:

\begin{pybox}
# Adiciona uma nova coluna chamada 'comments_ratio' ao DataFrame
socc['comments_ratio'] = None
\end{pybox}

Preenchemos a coluna com a proporção de comentários de nível superior sobre o total --- dividindo uma coluna pela outra:

\begin{pybox}
# Preenche a coluna 'comments_ratio' com a razao entre comentarios de
# nivel superior e o total de comentarios
socc['comments_ratio'] = socc['ntop_level_comments'] / socc['ncomments']
\end{pybox}

Uma amostra do resultado:

\begin{center}\small
\begin{tabular}{c c c c}
\toprule
\rowcolor{headrow}\thd{índice} & \thd{ntop\_level\_comments} & \thd{ncomments} & \thd{comments\_ratio} \\
\midrule
0 & 1378.0 & 2187.0 & 0.630087 \\
1 & 455.0 & 1103.0 & 0.412511 \\
2 & 433.0 & 1164.0 & 0.371993 \\
\bottomrule
\end{tabular}
\end{center}

Mas alguns artigos \textbf{não} receberam comentários --- nesses casos, dividiríamos zero por zero. Para essas linhas, @comments_ratio@ fica marcada como \textbf{@NaN@} (\textit{not a number}). O pandas \textbf{ignora automaticamente} os @NaN@ nos cálculos, como mostra o @describe()@:

\begin{pybox}
# Estatisticas descritivas da coluna 'comments_ratio'
socc['comments_ratio'].describe()
\end{pybox}
\begin{outbox}
count    7797.000000
mean        0.537057
std         0.205398
min         0.083333
25%         0.384615
50%         0.485714
75%         0.647059
max         1.000000
Name: comments_ratio, dtype: float64
\end{outbox}

Note o @count@: só 7\,797 itens (dos 10\,339) entraram no cálculo --- os @NaN@ foram ignorados.

\subsection{Aplicando processamento de linguagem a uma coluna}
E se quiséssemos fazer PLN e guardar os resultados? Selecionamos artigos no primeiro quartil de @comments_ratio@ \textbf{e} com mais de 200 comentários, usando @&@:

\begin{pybox}
# Filtra o DataFrame por artigos muito comentados e atribui o resultado a 'talk'
talk = socc.loc[(socc['comments_ratio'] <= 0.384) & (socc['ncomments'] >= 200)]
\end{pybox}

Importamos o spaCy e carregamos um modelo \textbf{médio} para o inglês:

\begin{pybox}
# Importa a biblioteca spaCy
import spacy

# Note que agora carregamos um modelo de tamanho medio!
nlp = spacy.load('en_core_web_md')
\end{pybox}

Ao criar uma coluna nova em @talk@, o pandas emite um \textbf{@SettingWithCopyWarning@}, porque @talk@ é apenas um recorte (\textit{slice}) do @DataFrame@ original. A solução é criar uma \textbf{cópia profunda} com @.copy()@:

\begin{pybox}
# Cria uma copia profunda do DataFrame
talk = talk.copy()

# Cria uma nova coluna chamada 'processed_title'
talk['processed_title'] = None
\end{pybox}

Para processar os títulos, usamos @apply()@, que \textbf{aplica} o que recebe a cada linha da coluna. Passamos o modelo @nlp@:

\begin{pybox}
# Aplica o modelo de linguagem 'nlp' ao conteudo da coluna 'title'
talk['processed_title'] = talk['title'].apply(nlp)
\end{pybox}

Cada célula da coluna @processed_title@ agora contém um objeto \textbf{@Doc@} do spaCy:

\begin{pybox}
# Pega o valor da coluna 'processed_title' na linha de indice 2
talk.at[2, 'processed_title']

# Verifica o tipo do objeto contido
type(talk.at[2, 'processed_title'])
\end{pybox}
\begin{outbox}
Too many first nations people live in a dream palace
spacy.tokens.doc.Doc
\end{outbox}

\subsection{Definindo a própria função e aplicando com \textnormal{@apply()@}}
Vamos definir uma função para buscar os \textbf{lemas de cada substantivo} do título. Funções são definidas com @def@, seguido do nome e dos parâmetros entre parênteses:

\begin{pybox}
# Define uma funcao 'get_nouns' que recebe um unico objeto como entrada.
# Referimo-nos a essa entrada pela variavel 'nlp_text'.
def get_nouns(nlp_text):

    # Primeiro garantimos que a entrada e do tipo correto,
    # usando 'assert' para checar o tipo
    assert type(nlp_text) == spacy.tokens.doc.Doc

    # Prepara uma lista vazia para os lemas
    lemmas = []

    # Percorre o objeto Doc
    for token in nlp_text:

        # Se a classe gramatical detalhada do token for substantivo (NN)
        if token.tag_ == 'NN':

            # Acrescenta o lema do token a lista de lemas
            lemmas.append(token.lemma_)

    # Ao fim do laco, retorna a lista de lemas
    return lemmas
\end{pybox}

Aplicamos com @apply()@, criando a coluna @nouns@ automaticamente pela atribuição:

\begin{pybox}
# Aplica a funcao 'get_nouns' a coluna 'processed_title'
talk['nouns'] = talk['processed_title'].apply(get_nouns)
\end{pybox}

Uma amostra dos títulos processados e seus substantivos:

\begin{center}\small
\begin{tabular}{c p{7.5cm} l}
\toprule
\rowcolor{headrow}\thd{índice} & \thd{title} & \thd{nouns} \\
\midrule
2 & Too many first nations people live in a dream\ldots & [dream, palace] \\
5 & Fifty years in Canada, and now I feel like a s\ldots & [class, citizen] \\
8 & A nation of \$100,000 firefighters & [nation] \\
\bottomrule
\end{tabular}
\end{center}

\begin{notebox}
Não é necessário criar uma coluna vazia antes --- o pandas cria a coluna nova automaticamente na atribuição, como em @talk['nouns']@.
\end{notebox}

\subsection{Extraindo dados para estruturas nativas do Python}
O método @tolist()@ extrai o conteúdo de uma @Series@ para uma \textbf{lista}:

\begin{pybox}
# Converte a pandas Series em uma lista
noun_list = talk['nouns'].tolist()

# Mostra os dez primeiros
noun_list[:10]
\end{pybox}
\begin{outbox}
[['dream', 'palace'], ['agency'], ['class', 'citizen'], ['right'],
 ['nation'], [], ['reform'], ['leader', 'parade'], ['pm'],
 ['government', 'monopoly']]
\end{outbox}

Temos uma \textbf{lista de listas}. Vamos ``achatar'' tudo em uma única lista @final_list@ com @extend()@:

\begin{pybox}
# Prepara a lista vazia
final_list = []

# Percorre cada lista dentro da lista de listas
for nlist in noun_list:

    # Estende a lista final com a lista atual
    final_list.extend(nlist)

# Verifica o tamanho da lista
len(final_list)
\end{pybox}
\begin{outbox}
884
\end{outbox}

Para plotar os 10 substantivos mais frequentes, convertemos em @Series@, contamos com @value_counts()@ e plotamos:

\begin{pybox}
# Converte a lista em uma pandas Series, conta os substantivos unicos com
# value_counts(), pega os 10 mais frequentes [:10] e plota em barras.
pd.Series(final_list).value_counts()[:10].plot(kind='bar')
\end{pybox}

% =========================================================== 4 ==============
\section{Salvando DataFrames}

@DataFrame@s podem ser salvos como objetos \textbf{serializados} (\textit{pickled}) com o método @to_pickle()@, que recebe o caminho do arquivo:

\begin{pybox}
# Grava o DataFrame em disco usando pickle
df.to_pickle('data/pickled_df.pkl')
\end{pybox}

\begin{warnbox}[Segurança (importante no CiberExt)]
O formato \textit{pickle} serializa objetos Python arbitrários e, ao ser \textbf{carregado}, pode \textbf{executar código arbitrário}. \textbf{Nunca} faça @read_pickle()@ (nem @pickle.load@, @joblib.load@, @numpy@ com @allow_pickle=True@ etc.) em arquivos de \textbf{fontes não confiáveis} --- um @.pkl@ malicioso compromete a máquina. Use \textit{pickle} só com dados que \textbf{você mesmo} gerou; para intercâmbio entre pessoas/sistemas, prefira formatos que não executam código, como \textbf{CSV}, \textbf{JSON} ou \textbf{Parquet}.
\end{warnbox}

Verificamos lendo de volta com @read_pickle()@:

\begin{pybox}
# Le o DataFrame serializado e atribui o resultado a 'df_2'
df_2 = pd.read_pickle('data/pickled_df.pkl')
\end{pybox}

Comparando os dois @DataFrame@s com @==@, obtemos um booleano por célula:

\begin{pybox}
# Compara os DataFrames 'df' e 'df_2'
df == df_2
\end{pybox}

\begin{center}\small
\begin{tabular}{c c c}
\toprule
\rowcolor{headrow}\thd{índice} & \thd{filename} & \thd{text} \\
\midrule
0 & True & True \\
1 & True & True \\
2 & True & True \\
\bottomrule
\end{tabular}
\end{center}

Tudo @True@ --- os dados foram salvos e recarregados com sucesso.

% ======================================================= resumo ============
\section*{Resumo da unidade}
\addcontentsline{toc}{section}{Resumo da unidade}
\begin{objbox}[Nesta unidade, você aprendeu a]
\begin{enumerate}[leftmargin=1.4em,itemsep=2pt,topsep=3pt]
  \item importar dados de um único arquivo (@read_csv()@) ou de vários (@Path@/@glob@ + laço com @.at@);
  \item examinar @DataFrame@s: @columns@, colunas como @Series@, @value_counts()@, @describe()@, @.plot()@ e filtragem com @.loc@ (@&@, @==@);
  \item estender: adicionar colunas, calcular entre colunas, lidar com @NaN@, usar @.copy()@ e aplicar funções com @apply()@;
  \item extrair com @tolist()@/@extend()@ e salvar/carregar com @to_pickle()@/@read_pickle()@ (com cuidado de segurança).
\end{enumerate}
\end{objbox}

\begin{notebox}[Exercícios sugeridos]
\begin{enumerate}[leftmargin=1.4em,itemsep=3pt,topsep=3pt]
  \item Carregue os três arquivos de notícias da Unidade 1 em um @DataFrame@ (colunas @filename@ e @text@) e conte quantos caracteres cada texto tem, numa nova coluna @n_chars@.
  \item No @DataFrame@ @socc@, use @.loc@ para selecionar os artigos de @GLOBE EDITORIAL@ com mais de 500 comentários.
  \item Escreva uma função @get_verbs()@ (análoga a @get_nouns@) que colete os lemas dos \textbf{verbos} (@token.pos_ == 'VERB'@) e aplique-a a uma coluna de títulos processados.
  \item Salve um @DataFrame@ em CSV com @to_csv('saida.csv', index=False)@ e recarregue com @read_csv()@, conferindo as colunas.
\end{enumerate}
\end{notebox}

\vspace{6pt}
\noindent\rule{\linewidth}{0.4pt}\\[4pt]
{\footnotesize\textbf{Referências.} \textit{Applied Language Technology} (MOOC), Universidade de Helsinque --- \href{https://applied-language-technology.mooc.fi/}{applied-language-technology.mooc.fi}. Kolhatkar, V. et al. (2020). \textit{The SFU Opinion and Comments Corpus.} Corpus Pragmatics. McKinney, W. et al. \textit{pandas: powerful Python data analysis toolkit} --- \href{https://pandas.pydata.org/}{pandas.pydata.org}.\par
Material produzido para o Programa CiberExt 26-29 / FEELT38103 --- Universidade Federal de Uberlândia.}

\end{document}
